\section{Linear exception counts from leading-degree separation}
\label{app:triangular-differential}

A first-order identity can force a linear exceptional-label bound even
when every agreement coordinate is singular. The criterion below separates
the nonlinear terms from the highest coefficient equations. Those equations
determine the nonconstant message coefficients with one denominator of
degree $O(D)$. The residual solution locus then has bounded degree in
one remaining message parameter and degree $O(D)$ in the challenge.
The conclusion concerns actual solutions of the identity: an application
to all nearby Reed--Solomon candidates must first establish that they all
satisfy such an equation.

\begin{theorem}[Leading-degree separation]\label{thm:triangular-newton}
Let $\F$ be any field, $D\ge1$, and $\ell\ge2$. Consider
\begin{equation}\label{eq:triangular-identity}
 T(X,z)P'+\sum_{j=0}^{\ell}A_j(X,z)P^j=0,
 \qquad \deg P\le D.
\end{equation}
Suppose $A_\ell\ne0$, $\deg_X T=t$ with leading coefficient
$\tau(z)\ne0$, and
\begin{equation}\label{eq:triangular-degree-separation}
 \deg_X A_1\le t-1,\qquad \deg_X A_0\le t+D-1,\qquad
 \deg_X A_j+jD<t\quad(2\le j\le\ell,\ A_j\ne0).
\end{equation}
Assume all coefficient polynomials have $z$-degree at most an integer $a\ge0$, and
that every pivot
\[
 d_i(z)=i\tau(z)+(A_1)_{t-1}(z),\qquad 1\le i\le D,
\]
is a nonzero polynomial. Put $M=a(\ell D+1)$ and $L=aD+1$.
On any $n$ distinct evaluation points, for any received line $f+zg$ and
integer $A>D$, outside a set of at most
\begin{equation}\label{eq:triangular-exception-bound}
 a(D+1)+2\ell M+\frac{n(M+\ell L)}{A-D}+\ell n
\end{equation}
labels, every solution of \eqref{eq:triangular-identity} with at least $A$
agreements admits degree-at-most-$D$ witnesses $F_0,G_0$ such that
$P=F_0+zG_0$ and its full agreement set is
\[
 \{x:f(x)=F_0(x),\ g(x)=G_0(x)\}.
\]
In particular the bound is $O_{a,\ell,\eta}(n)$ if
$A-D\ge\eta n$ and $D\le n$.
\end{theorem}
\begin{proof}
Write $P=c+\sum_{i=1}^D p_iX^i$. By
\eqref{eq:triangular-degree-separation}, the nonlinear terms have degree
less than $t$. Equating coefficients of $X^{t+i-1}$, in decreasing $i$,
gives a triangular linear system with pivot $d_i$. It does not involve
$c$. Every off-diagonal coefficient and inhomogeneous term has
$z$-degree at most $a$.

Set $\Delta_i=\prod_{h=i}^D d_h$ and $\Delta=\Delta_1$.
Descending induction gives
\[
 p_i=N_i/\Delta_i,\qquad \deg_zN_i\le a(D-i+1).
\]
Indeed, all higher denominators divide $\Delta_{i+1}$, so multiplying
the next equation by that product gives precisely the stated degree.
Passing to the single denominator $\Delta$ yields
\begin{equation}\label{eq:triangular-parametrization}
 P=c+\frac{W(X,z)}{\Delta(z)},\qquad W(0,z)=0,
 \qquad \deg_zW,\deg_z\Delta\le aD.
\end{equation}
Discard the at most $aD$ roots of $\Delta$. On the complement this is
an exact parametrization of the actual candidate locus by $(z,c)$.
No division by the integer $i$ is required: invertibility of the specified
pivot is the only characteristic condition needed.

After substituting \eqref{eq:triangular-parametrization} and multiplying
by $\Delta^\ell$, the remaining coefficient equations are the
coefficients in $X$ of
\[
 \Delta^{\ell-1}T W'
 +\sum_{j=0}^{\ell}A_j(c\Delta+W)^j\Delta^{\ell-j}.
\]
Each has bidegree at most $(M,\ell)$ in $(z,c)$.
Choose a nonzero coefficient $h(z)$ of $A_\ell$ in $X$ and also discard
its at most $a$ roots. The corresponding residual has leading
$c$-coefficient $h\Delta^\ell$, nonzero on this open set.

Work over an algebraic closure and take the gcd $U$ of all nonzero
residuals. Its factors give their common curve components. Since $U$
divides the chosen residual, its summed bidegrees are at most $(M,\ell)$.
Every vertical factor is supported on a root of $h\Delta$ and has already
been excluded. There are therefore at most $\ell$ reduced horizontal
components.

Divide the residuals by $U$. Their gcd is one. If one quotient is a
nonzero constant, there are no common zeros off $U$. Otherwise choose a
nonzero quotient and a generic constant linear combination of the other
quotients sharing none of its irreducible factors. Such a combination
exists because each factor is avoided by some quotient, and the
algebraic closure is infinite. Both polynomials have bidegrees at most
$(M,\ell)$. Their proper intersection in
$\mathbb P^1_z\times\mathbb P^1_c$ has total multiplicity at most
$2M\ell$. Every isolated point of the original locus lies off its curve
locus and in this intersection. Thus at most $2M\ell$ additional labels
come from isolated points. This argument removes common factors before
using intersection bounds; linear independence alone would not suffice.

At an evaluation coordinate $x$, agreement is the equation
\begin{equation}\label{eq:triangular-coordinate-agreement}
 \Delta c+W(x,z)-\Delta(f(x)+zg(x))=0,
\end{equation}
of bidegree at most $(L,1)$. For a curve component of bidegrees $(s,r)$,
unless agreement is identically true, its resultant with
\eqref{eq:triangular-coordinate-agreement} has degree at most $s+rL$.
Away from $\Delta=0$, each root gives at most one agreeing point because
the agreement equation is linear in $c$. Summing over components gives
at most $M+\ell L$ nonpersistent agreeing points per coordinate.

If a component has at most $D$ persistent agreement coordinates, each
$A$-near point on it supplies at least $A-D$ nonpersistent incidences.
Their labels number at most $n(M+\ell L)/(A-D)$. If it has more than $D$
persistent coordinates, interpolation forces its entire candidate
polynomial to be a base-field affine codeword graph $F_0+zG_0$.
Outside at most $n$ labels having an extra agreement, its full agreement
set is exactly the coefficientwise common agreement set of those
witnesses. There are at most $\ell$ components. Adding the discarded
labels, isolated points, and curve contributions proves
\eqref{eq:triangular-exception-bound}.
\end{proof}

\begin{corollary}[The unnormalized two-branch family]
\label{cor:triangular-two-branch}
Let $T$ be the monic squarefree locator of the $n$ evaluation points,
assume $n>4D$, and suppose the characteristic is zero or $p>D$.
Let $U_z,V_z\in\F[X]$ be affine in $z$, each of $X$-degree at most
$n-1$. For solutions of
\begin{equation}\label{eq:triangular-two-branch}
 TP'=(P-U_z)(P-V_z),\qquad\deg P\le D,
\end{equation}
on any chosen received line and threshold $A>D$, the full common-witness
property holds outside at most
\[
 18D+10+\frac{n(8D+4)}{A-D}+2n
\]
labels. Thus the family has only $O_\eta(n)$ exceptions when
$A-D\ge\eta n$.
\end{corollary}
\begin{proof}
Any nonzero coefficient of $U_zV_z$ above $X$-degree $n+D-1$ is a
polynomial in $z$ of degree at most two, and every solution label must
be its root. Otherwise the generic product degree is at most $n+D-1$.
If both generic branch degrees exceed $D$, each is at most $n-2$.
Apply Theorem~\ref{thm:triangular-newton} with
$A_2=-1$, $A_1=U_z+V_z$, $A_0=-U_zV_z$, $\ell=2$ and $a=2$.
The pivots are the nonzero constants $1,\ldots,D$. Substitution in
\eqref{eq:triangular-exception-bound} gives the displayed bound.

If a branch, say $U_z$, has degree at most $D$, it is an affine codeword
pencil. Any solution $P\ne U_z$ agrees with $U_z$ at at most $D$ domain
points and therefore with $V_z$ at at least $n-D$ points.
Such a solution is unique at each label, since two would agree at
$n-2D>D$ points. If there are at least three such labels, their
candidates satisfy the same affine relation as their labels on at least
$n-3D>D$ common agreement points with $V_z$. It is therefore a polynomial
identity. Every such candidate lies on the affine codeword pencil
determined by the first two; if there are at most two labels, discard them.
All solutions lie on at most two affine codeword pencils outside at most
two labels. For any received line, each pencil has the full-witness
property outside at most $n$ extra-agreement labels. This gives $2n+2$,
which is smaller than the displayed bound.
\end{proof}

\paragraph{Scope and the nonlinear boundary.}
The general theorem imposes no nonsingular-agreement requirement and no
separate characteristic restriction beyond its explicit nonzero pivots.
It does impose leading-degree separation, fixed challenge and value
degrees, and a derivative coefficient independent of $P$. Under $p>D$
there can be at most one identically zero pivot, since
$d_j-d_i=(j-i)\tau$; that resonant case is not included.
The quadratic isolated family of
Theorem~\ref{thm:quadratic-isolated-first-order} has derivative coefficient
$(z-X^2)R$ of degree $D+3$ and a $P^2$ coefficient of $X$-degree two.
Its nonlinear term can reach degree $2D+2$, so it violates the strict
separation condition. Its quadratic isolated count is not a contradiction.
The conclusion here is a family-specific upper bound for actual nearby
solutions, not a general linear first-order MCA theorem.

\paragraph{Checks.}
Independent symbolic elimination verifies the recurrence and residual
bidegrees through $D=4$, including nonmonic moving coefficients, value
degrees two through four, and persistent nonlinear components. An
independent finite-field check exhausts $29968$ polynomial--label pairs
and checks actual agreements, including persistent codeword pencils.
The scripts and reports are in \path{research/two_branch_recurrence/}.
The proofs above, rather than the finite fixtures, establish the
quantified statements.
